\scanpage{75}
\mspage{71}
De plus, l'on a \deleted{alors}, pour $A\subset E$
\[
\underset{\substack{\vrule height .45ex width .08ex\rule{7.2em}{.08ex}\vrule height .45ex width .08ex\\
\Vert\,\mathrm{def}\\
\operatorname{Supp}_{E'}^{\circ}\bigl(\bar{\alpha}(\operatorname{supp}_{E}^{\circ}(A)\bigr)}}{\alpha_{\unsure{\ast|\times}}\bigl(\operatorname{supp}_{E}^{\circ}(A)\bigr)}
=\deleted{\unsure{s}}\operatorname{Supp}_{E'}^{\circ}\bigl(\bar{\alpha}(\deleted{\unsure{E|A}}\unsure{A|E})\bigr)
\tag*{\makebox[0pt][r]{\makebox[\displaywidth][l]{\unsure{(xxx)|($\times\times\times$)}}}}
\]

\uline{NB} J'ignore si la validité de cette formule, pour une \deleted{$\alpha_{\ast}$} quasi-corr. disjonctive donnée $\alpha_{\ast}$, indique que c'est une corr. disjonctive

En analogie avec le cas des disjoncteurs triviaux, où on distingue le cas \unsure{d'}une correspondance qui est une application, ou \unsure{transposée|transposé} d'une application, je dirai que $\alpha_{\ast}$ est un \uline{morphisme} si $ {}^{t}\alpha_{\ast}=\beta_{\ast}$ \unsure{a} \deleted{\unsure{la}} les propriétés suivantes : il commute à $C$ (donc aux $\operatorname{Inf}$ quelconques)

% archive sheet number: 74
